\section{Infinitesimal study of schemes}
\source{fga-explained:page-0153}\source{fga-explained:page-0154}\source{fga-explained:page-0155}\source{fga-explained:page-0156}\source{fga-explained:page-0157}\source{fga-explained:page-0158}\source{fga-explained:page-0159}\source{fga-explained:page-0160}\source{fga-explained:page-0161}\source{fga-explained:page-0162}\source{fga-explained:page-0163}\source{fga-explained:page-0164}\source{fga-explained:page-0165}\source{fga-explained:page-0166}\source{fga-explained:page-0167}\source{fga-explained:page-0168}

\begin{convention}\label{fga-conv-6-0}\source{fga-explained:page-0153}\uses{}
Fix a field $k$ and write $\mathfrak m_A$ for the maximal ideal of a local
Artinian $k$-algebra $A$ with residue field $k$.
\end{convention}\begin{definition}[Artin deformation functor]\label{fga-def-6-1-2}\dcref{6.1.2}\source{fga-explained:page-0153}\uses{}
Let $(\operatorname{Art}/k)$ be the category of local Artinian $k$-algebras
with residue field $k$.  A deformation functor is a covariant functor
$D:(\operatorname{Art}/k)\to\operatorname{Sets}$ with $D(k)$ a singleton.
\end{definition}\begin{proposition}\label{fga-prop-6-1-4}\dcref{6.1.4}\source{fga-explained:page-0154}\uses{fga-def-6-1-2}
For a scheme $Y$ and $k$-point $y$, the functor of flat deformations of $y$
over local Artin rings is a deformation functor; liftings along a small
extension are the infinitesimal deformations of $y$.
\end{proposition}
\begin{definition}[tangent space]\label{fga-def-6-1-9}\dcref{6.1.9}\source{fga-explained:page-0155}\uses{fga-def-6-1-2}
The tangent space of $D$ is $t_D=D(k[\varepsilon]/(\varepsilon^2))$; it is a
$k$-vector space when $D$ is a deformation functor of schemes or modules.
\end{definition}
\begin{lemma}\label{fga-lem-6-1-17}\dcref{6.1.17}\source{fga-explained:page-0157}\uses{fga-def-6-1-9}
For a square-zero extension $0\to I\to A'\to A\to0$, the set of liftings of a
deformation, when nonempty, is a torsor under $t_D\otimes_k I$.
\end{lemma}
\begin{theorem}\label{fga-thm-6-1-19}\dcref{6.1.19}\source{fga-explained:page-0158}\uses{fga-lem-6-1-17,fga-def-5-1}
For a closed subscheme $Z\subset X$ over $k$, the tangent space to the Hilbert
functor at $[Z]$ is $H^0(Z,\mathcal N_{Z/X})=\operatorname{Hom}(\mathcal I_Z,
\OO_Z)$; obstructions lie in $H^1(Z,\mathcal N_{Z/X})$.
\end{theorem}
\begin{proof}
An infinitesimal deformation is an ideal in $\OO_X\otimes k[\varepsilon]$
reducing to $\mathcal I_Z$.  Its first-order displacement is a map
$\mathcal I_Z\to\OO_Z$; changing local liftings changes this map by a Cech
 coboundary, and the obstruction to gluing lies in $H^1(\mathcal N_{Z/X})$.
\end{proof}
\begin{definition}[obstruction theory]\label{fga-def-6-1-21}\dcref{6.1.21}\source{fga-explained:page-0159}\uses{fga-lem-6-1-17}
An obstruction theory assigns to each small extension an obstruction class in a
vector space such that the class vanishes exactly when a lifting exists, and
then the liftings form a torsor under the tangent space tensored with the kernel.
\end{definition}

\section{Pro-representable functors}
\begin{definition}[pro-representability]\label{fga-def-6-2-1}\dcref{6.2.1}\source{fga-explained:page-0161}\uses{fga-def-6-1-2}
A deformation functor $D$ is pro-representable by a complete local noetherian
$k$-algebra $R$ if $D(A)\simeq\operatorname{Hom}_{\mathrm{cont},k}(R,A)$
naturally for every $A\in(\operatorname{Art}/k)$.
\end{definition}
\begin{theorem}[Schlessinger criterion]\label{fga-thm-6-2-4}\dcref{6.2.4}\source{fga-explained:page-0162}\uses{fga-def-6-2-1}
If $D$ preserves fiber products over small extensions, has finite-dimensional
tangent space, and satisfies the Schlessinger surjectivity condition, then $D$
is pro-representable (with uniqueness under the corresponding injectivity
condition).
\end{theorem}
\begin{proof}
Choose a basis of the tangent space and inductively construct a universal object
over quotients of a power-series ring.  The fiber-product condition lifts the
universal element across each small extension; completeness gives the inverse
limit.  The Schlessinger conditions prove existence and uniqueness of the
resulting natural map.
\end{proof}
\begin{corollary}\label{fga-cor-6-2-5}\dcref{6.2.5}\source{fga-explained:page-0163}\uses{fga-thm-6-2-4}
The deformation functor of a projective scheme with finite-dimensional tangent
and obstruction spaces is pro-representable whenever its automorphism condition
is finite.
\end{corollary}
\begin{corollary}\label{fga-cor-6-2-6}\dcref{6.2.6}\source{fga-explained:page-0163}\uses{fga-thm-6-2-4}
The completed local ring of a fine moduli space pro-represents the corresponding
deformation functor.
\end{corollary}

\section{Non-pro-representable functors}
\begin{definition}\label{fga-def-6-3-1}\dcref{6.3.1}\source{fga-explained:page-0164}\uses{fga-def-6-2-1}
A deformation functor is hullable if it has a smooth map from a pro-representable
functor inducing an isomorphism on tangent spaces; it is not required to be
universal on higher Artin rings.
\end{definition}
\begin{theorem}\label{fga-thm-6-3-4}\dcref{6.3.4}\source{fga-explained:page-0165}\uses{fga-def-6-3-1,fga-thm-6-2-4}
Schlessinger's conditions (H1)--(H3) give a hull, while adding (H4) gives a
pro-representing object.
\end{theorem}
\begin{corollary}\label{fga-cor-6-3-5}\dcref{6.3.5}\source{fga-explained:page-0165}\uses{fga-thm-6-3-4}
Automorphisms can prevent pro-representability even when a hull exists; the
obstruction is measured by the failure of uniqueness in (H4).
\end{corollary}

\section{Examples of tangent-obstruction theories}
\begin{definition}\label{fga-def-6-4-2}\dcref{6.4.2}\source{fga-explained:page-0166}\uses{}
For a morphism $X\to S$ and a square-zero extension by a quasi-coherent module
$J$, the cotangent complex $L_{X/S}$ controls liftings; the obstruction lies in
$\operatorname{Ext}^2(L_{X/S},J)$, liftings form a torsor under
$\operatorname{Ext}^1(L_{X/S},J)$, and automorphisms are
$\operatorname{Ext}^0(L_{X/S},J)$.
\end{definition}\begin{theorem}\label{fga-thm-6-4-5}\dcref{6.4.5}\source{fga-explained:page-0167}\uses{fga-def-6-4-2}
For a smooth morphism the cotangent complex is the sheaf of differentials in
degree zero, so liftings across square-zero thickenings are unobstructed and
form a torsor under $\operatorname{Hom}(\Omega_{X/S},J)$.
\end{theorem}
\begin{proposition}\label{fga-prop-6-4-7}\dcref{6.4.7}\source{fga-explained:page-0167}\uses{fga-def-6-4-2}
For a local complete intersection, the cotangent complex has amplitude
$[-1,0]$, so obstructions lie in $\operatorname{Ext}^1$ of the conormal
complex and there are no higher obstructions.
\end{proposition}
\begin{definition}\label{fga-def-6-4-8}\dcref{6.4.8}\source{fga-explained:page-0168}\uses{fga-def-6-1-21}
A tangent-obstruction theory on a functor consists of a tangent space, an
obstruction space, and functorial obstruction and torsor maps for small
extensions satisfying additivity under composition.
\end{definition}
\begin{theorem}\label{fga-thm-6-4-9}\dcref{6.4.9}\source{fga-explained:page-0168}\uses{fga-def-6-4-8,fga-def-5-1}
The Hilbert functor, the deformation functor of a coherent sheaf, and the
deformation functor of a morphism admit tangent-obstruction theories given by
the corresponding $\operatorname{Ext}$ groups of normal, endomorphism, and
cotangent complexes.
\end{theorem}
